\section{A triangular path gives exponentially many distinct costs}
\label{sec:cycles}

We construct a column code with $\exp(\Omega(m))$ pointwise ordered
scalar enumerators on $m$ logical coordinates. All path increments
belong to $\{2,4,6\}$. The argument uses the analytic cost expansion
and an exact triangular transfer matrix, and requires no arithmetic
independence assumption about optimized tanh values.

\begin{lemma}[A positive inverse Laplace kernel near zero]
\label{lem:laplace-kernel}
Suppose a real-on-the-real-axis function $w$ is holomorphic on a right
half-plane and satisfies
\[
 w(z)=Az+B+C/z+E(z),\qquad C>0,\qquad |E(z)|\le K/|z|^2
\]
uniformly there. There is a real continuous function $\rho$ of at most
exponential growth such that, for sufficiently large real $n$,
\[
 w(n)=An+B+\int_0^\infty e^{-nt}\rho(t)\,dt.
\]
Moreover $\rho(t)\ge C/2$ for all sufficiently small positive $t$.
\end{lemma}
\begin{proof}
For an admissible vertical line $\Re z=c$, put
\[
 r(t)=\frac1{2\pi}\int_{\mathbb R}E(c+iy)e^{(c+iy)t}\,dy.
\]
The integral is absolutely convergent. Closing rectangles proves
independence of $c$ for $t>0$. Fubini followed by Cauchy's formula gives
$E(z)=\int_0^\infty e^{-zt}r(t)\,dt$ for $\Re z>c$.
Conjugation symmetry makes $r$ real, and
\[
 |r(t)|\le\frac{K e^{ct}}{2c}.
\]
For small $t$ choose $c=1/t$, obtaining $|r(t)|\le eKt/2$.
Thus $\rho=C+r$ is positive near zero. A fixed admissible $c$ gives
the claimed exponential-growth bound and continuity follows by
dominated convergence.
\end{proof}

The interpolation in Theorem~\ref{thm:curvature} satisfies this lemma
after restricting to a sufficiently far-right half-plane.

\subsection{An exactly triangular path transfer}

Take a path on $V=2k+1$ vertices, where $k\ge2$. Give every loop
increment four, and give its successive edges increments
$6,2,6,2,\ldots,6,2$. For $t>0$ put
\[
 a=\tanh(2t),\qquad e_j=\tanh(\delta_jt/2),\qquad
 M(a,e)=\begin{pmatrix}1&ae\\a&e\end{pmatrix},\qquad
 J=\begin{pmatrix}0&1\\1&0\end{pmatrix}.
\]
If $B$ is the path incidence matrix, the syndrome enumerator is
\[
 Z_v(t)=\sum_{y\in\mathbb F_2^{2k}}
             \left(\prod_{j=1}^{2k}e_j^{y_j}\right)a^{|v+By|}.
\]
With $c_0=(1,a)^T$ and $c_1=(a,1)^T$, successive edge summation gives
\begin{equation}\label{eq:path-transfer-pair}
 \binom{Z_v}{Z_{v+e_1}}
 =J^{v_1}M(a,e_1)\cdots J^{v_{2k}}M(a,e_{2k})c_{v_{2k+1}}.
\end{equation}
The first coordinate here is the scalar enumerator whose order is
needed below.

Let
\[
 \mathsf H=\begin{pmatrix}1&1\\1&-1\end{pmatrix},\qquad
 \mathsf D=\begin{pmatrix}1&0\\0&-1\end{pmatrix}.
\]
Thus $\mathsf H^{-1}=\mathsf H/2$ and
$\mathsf HJ\mathsf H^{-1}=\mathsf D$. Direct calculation gives
\begin{equation}\label{eq:path-hadamard}
 \mathsf HM(a,e)\mathsf H^{-1}
 =\alpha\begin{pmatrix}1&x\\yx&y\end{pmatrix},\qquad
 \alpha=\frac{(1+a)(1+e)}2,\quad
 x=\frac{1-e}{1+e},\quad y=\frac{1-a}{1+a}.
\end{equation}
Here $y=e^{-4t}$ and $x=e^{-\delta_jt}$.

Fix every even-position syndrome bit to one. Write $x=e^{-2t}$.
The two steps of each block, with their positive factors $\alpha$
temporarily omitted, multiply to
\begin{equation}\label{eq:path-triangular}
 \begin{pmatrix}1&x^3\\x^5&x^2\end{pmatrix}
 \mathsf D
 \begin{pmatrix}1&x\\x^3&x^2\end{pmatrix}
 =\begin{pmatrix}1-x^6&x(1-x^4)\\0&-x^4(1-x^2)\end{pmatrix}.
\end{equation}
Consequently, if the preceding odd-position bit is $\eta_i$, the
complete block equals a common positive scalar $L(t)$ times
\[
 T_i(t)=\begin{pmatrix}1&b(t)\\0&\epsilon_i r(t)\end{pmatrix},
 \qquad \epsilon_i=(-1)^{\eta_i+1},
\]
where
\begin{equation}\label{eq:path-ratio}
 b(t)=\frac{x(1+x^2)}{1+x^2+x^4}>0,\qquad
 r(t)=\frac{x^4}{1+x^2+x^4}\in(0,1/3),
\end{equation}
and
\[
 L(t)=\alpha_6(t)\alpha_2(t)(1-x^6)>0.
\]
The continuous limits at zero are $b(0)=2/3$, $r(0)=1/3$, and
$L(t)=3t+O(t^2)$.

Fix the terminal syndrome bit to zero. Let the $k$ odd-position bits
$\eta_1,\ldots,\eta_k$ vary subject to one prescribed total parity.
Put
\[
 d_0=1,\qquad
 d_j=\prod_{i=k-j+1}^k\epsilon_i\quad(1\le j\le k-1),\qquad
 E=\prod_{i=1}^k\epsilon_i.
\]
The sign $E$ is common to the entire family, while
$(d_1,\ldots,d_{k-1})$ ranges freely over $\{-1,1\}^{k-1}$.
Multiplying the triangular matrices, then using
$\mathsf Hc_0=(1+a,1-a)^T$ and returning to the original first
coordinate in \eqref{eq:path-transfer-pair}, gives the exact scalar
formula
\begin{equation}\label{eq:path-scalar-digits}
 Z_\eta(t)=\frac{L(t)^k}{2}
 \left[(1+a)+(1-a)b(t)\sum_{j=0}^{k-1}d_jr(t)^j
                     +(1-a)E r(t)^k\right].
\end{equation}
The fixed total parity is what makes the final boundary term common.

\begin{lemma}[Pointwise strict scalar order]
\label{lem:path-pointwise-order}
The $2^{k-1}$ scalar enumerators in
\eqref{eq:path-scalar-digits} have the same strict lexicographic order
for every $t>0$. For one fixed $t_0>0$ and an absolute $C_1$, every
distinct pair satisfies
\begin{equation}\label{eq:path-pointwise-gap}
 |Z_\eta(t)-Z_{\eta'}(t)|\ge e^{-C_1k}t^k
 \qquad(0<t\le t_0).
\end{equation}
\end{lemma}
\begin{proof}
If $j\ge1$ is the first differing digit and $d_j=1$, $d'_j=-1$,
then $0<r\le1/3$ implies
\[
 \sum_l(d_l-d'_l)r^l
 \ge2r^j-2\sum_{l>j}r^l\ge r^j>0.
\]
All factors multiplying this sum in \eqref{eq:path-scalar-digits}
are positive. This proves the fixed pointwise order. On a sufficiently
small fixed interval, $b$, $1-a$, and $r$ are bounded below by positive
absolute constants and $L(t)\ge ct$. Since $j\le k-1$, these bounds
give \eqref{eq:path-pointwise-gap}.
\end{proof}

\subsection{Embedding into rectangle costs}

Add $P=2V+1=4k+3$ independent prefix coordinates and two distinguished
coordinates. The logical dimension is
\[
 m=V+P+2=6k+6.
\]
Use one effective loop column for every logical coordinate and all
$V-1=2k$ path-edge columns. The number of effective columns is
$\ell=m+V-1=8k+6$. Put
\[
 H=2\lceil m\log(m+2)\rceil,\qquad
 R=2\lceil K\log^2(m+2)\rceil,
\]
where $K$ is a sufficiently large fixed constant. Give the distinguished
loops effective increment $H$, every other loop increment four, and
the path edges alternating increments six and two. Realize each
distinguished loop by $g=\lceil H/R\rceil$ repeated physical columns,
whose positive even increments sum to $H$ and individually are at
most $R$. Thus the physical pair count is
\[
 p=\ell-2+2g=\Theta(m),\qquad
 n_0=pR+1=\Theta(m\log^2m).
\]
The shared carrier and common baseline $R$ give
\[
 n_s=n_0+H(s_1+s_2)
       +\sum_{j\text{ other}}\delta_j(a_j\cdot s),
 \qquad n_0=pR+1.
\]
The bits in this degree formula are interpreted as zero or one.
Splitting a distinguished column into repeated copies does not change
its displayed effective degree contribution. Every physical increment
is at most $R$, so the exact atom theorem applies for large enough
$K$. Also $H/n_0=O(1/\log m)$.

Let $u=e_1+e_2$. Each candidate $v_\eta$ has both distinguished bits
and every prefix bit one. Its path syndrome has even-position bits
one, terminal bit zero, and odd-position bits $\eta$ of the prescribed
total parity. Put
\[
 W(v)=2^{-m}\sum_s(-1)^{s\cdot v}w(n_s),\qquad
 \nu_\eta=W(v_\eta)-W(v_\eta+u).
\]
Neither candidate nor its translate by $u$ is zero or a column, so the
affine part $An+B$ has no coefficient there. Factoring the Laplace
integrand gives
\begin{equation}\label{eq:cycle-cost}
 \nu_\eta=-\int_0^\infty e^{-(n_0+H)t}\rho(t)
       A_{\rm prefix}(t)A_{\rm path}(t)Z_\eta(t)\,dt.
\end{equation}
Here $A_{\rm prefix}=[(1-e^{-4t})/2]^P$, and $A_{\rm path}$ is the
product of the common positive Fourier factors
$(1+e^{-\delta_jt})/2$ for path loops and edges. The two distinguished
factors give exactly
\[
 [(1-e^{-Ht})/2]^2-[(1+e^{-Ht})/2]^2=-e^{-Ht}.
\]

\begin{theorem}[Exponentially many ordered tanh rectangle costs]
\label{thm:cycle-order}
Assume $H\le C_2n_0$ for a fixed constant $C_2$ and
\begin{equation}\label{eq:path-baseline-growth}
 \frac{n_0}{m\log n_0}\longrightarrow\infty.
\end{equation}
For all sufficiently large $k$, all $2^{k-1}$ costs in
\eqref{eq:cycle-cost} are negative and pairwise distinct in a fixed
strict order. Thus there are $\exp(\Omega(m))$ competing cost levels.
\end{theorem}
\begin{proof}
The sign of $Z_\eta-Z_{\eta'}$ is fixed at every positive time.
On $t\in[1/n_0,2/n_0]$, the kernel satisfies $\rho\ge C/2$,
$A_{\rm prefix}\ge t^P$, and $A_{\rm path}$ is bounded below by a
positive absolute constant. The last assertion follows because its
logarithm is $O(mt)$ and all path increments are at most six.
Also $e^{-(n_0+H)t}$ is bounded below. Hence the positive near-zero
contribution to a cost difference has magnitude at least
\begin{equation}\label{eq:cycle-gap}
 e^{-C_3k}n_0^{-(P+k+1)}=\exp[-O(m\log n_0)].
\end{equation}

Choose fixed $t_0>0$ with $\rho\ge C/2$ on $(0,t_0]$.
Every Fourier coefficient of
$\exp[-t\sum_j\delta_j(a_j\cdot s)]$ is nonnegative and at most
one, because its factor expansion has nonnegative coefficients with
total sum one. Thus a difference of two complete coefficient
integrands is at most two in absolute value after the common
$e^{-n_0t}$ factor is removed. The exponential-growth bound on
$\rho$ bounds the adverse tail $t\ge t_0$ by
$\exp[-\Omega(n_0)]$. By \eqref{eq:path-baseline-growth}, this is
smaller than \eqref{eq:cycle-gap}. The integral therefore preserves
the fixed scalar order, with its overall negative sign.

The same argument without subtraction proves every $\nu_\eta<0$:
for a direct uniform lower bound, the all-loop term in $Z_\eta$ is
$a^{|v_{\rm path}|}\ge t^{2k+1}$ on the small interval.
The same lower bound with the two positive distinguished factors
proves $W(v_\eta)>0$. Consequently
$W(v_\eta+u)>W(v_\eta)\ge0$ throughout the affine family.
\end{proof}

\subsection{Calibration and resource compatibility}

The family of all candidates is itself an affine coset $U$: both
distinguished bits and all prefix bits are fixed to one, even path
bits are fixed to one, the terminal path bit is fixed to zero, and
the free odd path bits have one prescribed total parity. A loop at
a fixed bit moves outside $U$; a loop at a free bit changes its
prescribed parity. Every path edge changes a fixed even-position
bit, including the final edge. Thus $U\cup(U+u)\cup\{u\}$ is
independent in the Cayley graph of all physical columns. The prefix
separates $u$ from the two cosets, and no column equals $u$.
Repeated distinguished columns preserve this property.

The common-dual and physical-rounding arguments therefore apply
unchanged once their dominant-Walsh inequalities hold. The $P$
compulsory prefix factors give the uniform tail estimate
\[
 |\nu_v|\le O(2^P P!/n_0^{P+1})+\exp[-\Omega(n_0)],
 \qquad v\in U.
\]
When $n_0/m\to\infty$, this absorbs $2^m$ times every fixed
polynomial stability and rounding factor under
\eqref{eq:path-baseline-growth}.

For the displayed parameter choice, the effective increment cube has
$q=\ell=O(m)$ coordinates, total ordinary increment $S=O(m)$, and
total increment $D_\delta=2H+S=O(m\log m)=o(n_0)$. In particular
$q\log(n_0+D_\delta)=o(n_0)$. The inverse-Laplace Walsh estimates
therefore apply, with relative off-pair mass
\[
 O(S/H+H/n_0)=O(1/\log m).
\]
Higher-order coefficients push forward as sums under the column map;
negative singleton contributions occur only at individual columns.
Neither $u$ nor a live candidate is a column, so the common-dual
dominance and sign hypotheses hold for all large $m$.

The physical input dimension is
\[
 d_m=1+2pR+2H+S=\Theta(m\log^2m).
\]
Every target parity has degree $O(m\log^2m)$. The same threshold-step
tanh teacher consequently has width
$C_m=O(m^2\log^2m)$, coefficients bounded by that displayed width,
and uniform margin one. The designated mistake count is
\[
 T_m=2^{k-1}=2^{m/6-2}
 =\exp\!\left(\Omega\!\left(\frac{\sqrt{C_m}}{\log C_m}\right)\right).
\]
The target's bounded-fan-in Boolean circuit has $O(d_m)$ gates;
repeated columns merely contribute additional rail parities to the
same two distinguished logical bits.
